\color{unidarkgreen}{\subsection{Фиксация положения выключателей (ФПВ)}}
\color{black}

\subsubsection{Общие сведения}

\begin{enumerate}[label=\arabic{section}.\arabic{subsection}.\arabic{subsubsection}.\arabic*,  align=left, labelsep=4pt, leftmargin=0pt, itemindent=57pt]

\item
Логика ФПВ В1(ЛВ)ВН, ФПВ В2(ОВ)ВН, ФПВ В3, ФПВ В НН2(СН),  аналогичны, рассмотрим на
примере ФПВ В1(ЛВ)ВН.

\item
Ввод в работу ФПВ В1(ЛВ)ВН осуществляется с помощью уставки \textbf{«ФПВ В1(ЛВ)ВН»}. Оперативный
ввод/вывод ФПВ производится с помощью БУ (блоков управления функциями) «Вывод ФПВ В1(ЛВ)
ВН»». Предусматривается возможность формирования положения выключателей с контролем
разъединителей, которая задается уставкой \textbf{«Контроль Р В1(ЛВ)ВН»}. С помощью уставки \textbf{«Схема
Р В1(ЛВ) ВН»} имеется возможность выбрать схему выключателя, два или три разъединителя. С
помощью уставки \textbf{«Тип В1(ЛВ) ВН»} имеется возможность выбрать тип выключателя, линейный или
обходной. Так же от блок-контактов В1(ЛВ) ВН формируется внутренний пуск ЗНФ, который пофазно
контролирует расхождение контактов полюсов выключателей, возникающее при трехфазном включении
или отключении выключателя.

\begin{figure}[H]
\centering
\includegraphics[width=0.6\textwidth,height=0.6\textheight,keepaspectratio]{\fbpath/761.3471 ФПВ/img/FPV_1.pdf}
%\captionof{figure}{Функциональная схема ФПВ В1}\label{pic:FPV_1}
\end{figure}

\begin{figure}[H]
\centering

\includegraphics[width=0.6\textwidth,height=0.6\textheight,keepaspectratio]{\fbpath/761.3471 ФПВ/img/FPV_2.pdf}
\captionof{figure}{Функциональная схема ФПВ В1}\label{pic:FPV_2}
\end{figure}

\small
\begin{longtable}{|>{\centering\arraybackslash}m{5.3cm}|>{\centering\arraybackslash}m{4.2cm}|>{\centering\arraybackslash}m{1cm}|>{\centering\arraybackslash}m{3.3cm}|>{\centering\arraybackslash}m{1.8cm}|}
% Первая страница: полный заголовок
\caption{Перечень уставочных параметров ФПВ\hfill\vspace{-0.5\baselineskip}}\label{desc:tbl1}\\
\hline
\rowcolor{gray!30}
Наименование уставки & Диапазон & Шаг & Значение по умолчанию & Единица измерения \\
\hline
\endfirsthead

% Промежуточные страницы: «Продолжение таблицы»
\caption*{\hspace{3pt}\emph{Продолжение таблицы \ref{desc:tbl1}}\hfill\ }\vspace{-0.5\baselineskip} \\
\hline
\rowcolor{gray!30}
Наименование уставки & Диапазон & Шаг & Значение по умолчанию & Единица измерения \\
\hline
\endhead

%===m>ДТЗ|-
%===m
\end{longtable}
\normalsize	

\end{enumerate}